\section{Introduction}
\label{sec:introduction}

Large Language Models (LLMs) have undergone remarkable advancements in recent years, evolving from text generators to systems capable of interacting with their environment through API calls and function execution. 
This capability, commonly referred to as ``tool use'' \citep{schick2023toolformer, qin2024tooll, mialon2023augmented}, enables LLMs to perform actions ranging from simple information retrieval to complex operations that modify the state of the systems they interact with such as performing monetary transactions or modifying databases.
Major language models \citep{achiam2023gpt, grattafiori2024llama, jiang2024mixtral, team2024gemini} 
now allow their users to integrate them with external tools, databases, and services through their \emph{Function-Calling} (FC) interface \citep{openai2025fc,gemini2025fc}. Such interfaces have become the most widespread approach for accessing the LLM's tool-use capabilities \citep{carrigan2024tool}.
The integration of tool-use capabilities represents a leap toward autonomous AI systems that can take actions on behalf of users. 

However, with this increased autonomy come increased risks and responsibilities. 
Unlike errors in text generation that may be less consequential, errors in tool use can result in irreversible changes, leading to data corruption, financial losses, or safety incidents \citep{andriushchenko2024, aichberger2025attacking}. In this context, Uncertainty Quantification (UQ) emerges as a critical component for responsible deployment of LLM-based systems with tool-use capabilities. 
UQ methods aim to estimate a model's confidence in its decisions,
enabling systems to abstain from actions or defer to human oversight when uncertainty is high. 
Effective UQ methods can serve as guardrails, preventing models from executing potentially harmful actions when they lack sufficient confidence. 

Despite the clear importance of UQ in FC scenarios, there is a notable absence of systematic evaluations for assessing UQ methods in this specific context, as well as of UQ methods developed specifically for it. 
While substantial research has addressed uncertainty estimation for classification and question-answering tasks \citep{kirchhof2024pretrained, kuhn2023semantic, ye2024benchmarking}, the tool-use setting introduces unique challenges. 
FC requires structured outputs that adhere to predefined syntax, often involving complex reasoning about parameter selection and the appropriateness of specific API calls. 
It is therefore unclear how existing UQ methods fare in this setting.

In this paper we aim to address this gap. We  conduct, to our knowledge, the first comprehensive evaluation of uncertainty quantification methods in LLM tool-use scenarios. 
To do so, we create a benchmark based upon the Berkeley Function Calling Leaderboard dataset (BFCL) \citep{bfcl}, which provides a collection of FC examples, mostly in Python.
Our benchmark evaluates various UQ approaches based on how well their confidence estimates predict the correctness of function calls outputted by LLMs.
We reimplement several UQ methods, ranging from simple single-sample approaches based on model logits to more sophisticated multi-sample approaches that incorporate semantic equivalence between samples, and test how effectively they discriminate correct from incorrect function calls in a variety of setups.

Our findings indicate that in the FC setting, unlike in natural language Q\&A settings, the more involved entropy-based multi-sample methods we evaluate (predictive and semantic entropy) do not outperform simple single-sample UQ methods, such as the Greedy Negative Log-Likelihood (G-NLL).
We then investigate if the performance of existing UQ methods in the FC setting can be increased by tailoring them to the characteristics of FC outputs. By adapting the clustering strategy in Semantic Entropy and the token selection in G-NLL, we achieve performance improvements for both multi-sample and single-sample UQ methods.



\noindent In summary, our contributions are as follows:
\begin{enumerate}[leftmargin=*, topsep=0pt, itemsep=2pt, partopsep=0pt, parsep=0pt]
    \item We build the first benchmark for UQ methods in Function-Calling settings. 
    \item We find that
    Semantic Entropy (SE), the best-performing of the entropy-based multi-sample UQ methods we evaluate, offers no clear advantage over G-NLL as the best single-sample method.
    \item We analyze the properties of FC outputs
    and use them to devise variants of existing UQ methods that are tailored to the FC setting and improve their performance in this setting. Namely, we adapt the semantic-entailment method of SE to the FC setting by using Abstract Syntax Trees (AST) to cluster FC outputs, and adjust G-NLL and other single-sample UQ methods to the FC setting by defining semantically meaningful tokens and including only those in  calculating UQ scores. AST-based clustering consistently improves the multi-sample methods, while restricting single-sample scores to semantically meaningful tokens leaves their ranking performance essentially unchanged and instead improves their calibration and, to a smaller extent, their utility under abstention.
\end{enumerate}

\textbf{Artifact availability.} We release the benchmark harness (the BFCL-derived task splits and our selective-prediction evaluation), our implementations of all UQ methods evaluated here (including the AST-based clustering and the semantically meaningful token extractors for both output formats): \texttt{[REPOSITORY URL WILL BE INSERTED FOR CAMERA READY]}.

